\medskip

L'espace est muni d'un repère orthonormé \Oijk.

On considère les trois points A(3~;~0~;~0), B(0~;~2~;~0) et C(0~;~0~;~2).





















\bigskip

\textbf{Partie 1 : Distance du point O au plan (ABC)}

\medskip

\begin{enumerate}
\item Soit le vecteur $\vect{n}(2~;~3~;~3)$.

\begin{list}{\textbullet}{}
\item Le vecteur $\vectt{AB}$ a pour coordonnées $\left ( -3\,;\, 2 \,;\, 0\right )$.

$\vectt{AB}\cdot\vect{n}=-3 \times 2 + 2\times 3 + 0 \times 3 = 0$ donc $\vectt{AB} \perp \vect{n}$
\item Le vecteur $\vectt{AC}$ a pour coordonnées $\left ( -3\,;\, 0 \,;\, 2\right )$.

$\vectt{AC}\cdot\vect{n}=-3 \times 2 + 0\times 3 + 2 \times 3 = 0$ donc $\vectt{AC} \perp \vect{n}$
\end{list}

Le vecteur $\vect{n}$ est donc orthogonal aux deux vecteurs $\vectt{AB}$ et $\vectt{AC}$ non colinéaires, donc il est normal au plan (ABC).

\item
Le plan (ABC) est l'ensemble des points M de coordonnées $\left ( x\,;\, y \,;\, z\right )$ tels que $\vectt{AM} \perp \vect{n}$.

Le vecteur $\vectt{AM}$ a pour coordonnées $\left ( x-3\,;\, y \,;\, z\right )$.

$\vectt{AM} \perp \vect{n}
\iff \vectt{AM} \cdot \vect{n} = 0
\iff 2\left (x-3\right ) + 3y + 3z=0
\iff 2x +3y +3z -6=0$

Le plan (ABC) a donc pour équation cartésienne:

\[M(x~;~y~;~z) \in \text{ABC} \iff 2x +3y +3z -6  = 0.\]

\item La droite $d$ passant par O et de vecteur directeur $\vect{n}$ est l'ensemble des points M de coordonnées $\left ( x\,;\, y \,;\, z\right )$ tels que $\vectt{OM}$ et $\vect{n}$ soient colinéaires, c'est-à-dire tels que \\
$\vectt{OM}=t.\vect{n}$ où $t\in\R$.

Le vecteur $\vectt{OM}$ a pour coordonnées $\left ( x\,;\, y \,;\, z\right )$.
Or
$\vectt{OM}=t.\vect{n}
\iff 
\left \{
\begin{array}{l !{=} l}
x & 2t\\
y & 3t\\
z & 3t
\end{array}
\right . $

La droite $d$ passant par O et de vecteur directeur $\vect{n}$ a donc pour représentation paramétrique:
$\left \{
\begin{array}{l !{=} l}
x & 2t\\
y & 3t\\
z & 3t
\end{array}
\right . \quad t\in\R$

\item On note H le point d'intersection de la droite $d$ et du plan (ABC).

Les coordonnées de H vérifient le système
$\left \{
\begin{array}{r !{=} l}
x & 2t\\
y & 3t\\
z & 3t\\
2x+3y+3z-6& 0
\end{array}
\right . $

On a donc: $2\left (2t\right )+3\left (3t\right )+3\left (3t\right )-6= 0$, donc 
$22t=6$ donc $t=\dfrac{3}{11}$.

Les coordonnées de H sont donc: $\left (\dfrac{6}{11}\,;\, \dfrac{9}{11}\,;\, \dfrac{9}{11} \right )$.


\item La distance du point O au plan (ABC) est OH.

$\text{OH}^2=\left (\dfrac{6}{11} \right )^2 + \left (\dfrac{9}{11} \right )^2 + \left (\dfrac{9}{11} \right )^2 = \dfrac{36+81+81}{11^2}=\dfrac{198}{11^2}$
donc $\text{OH} = \dfrac{\sqrt{198}}{11}=\dfrac{3\sqrt{22}}{11}$.
\end{enumerate}

\bigskip

\textbf{Partie 2 : Démonstration de la propriété}



\begin{enumerate}
\item
En prenant pour base le triangle OAB et pour hauteur OC, le volume du tétraèdre OABC vaut $\dfrac{\text{OC}\times \text{aire(OAB)}}{3}$.

$\text{OC}=2$ et $\text{aire(OAB)} = \dfrac{\text{OA} \times \text{OB}}{2}=\dfrac{3\times 2}{3}=3$

Le volume du tétraèdre OABC  est donc égal à $\dfrac{2\times 3}{3}=2$.

\item 
En prenant pour base le triangle ABC et pour hauteur OH, le volume du tétraèdre OABC vaut $\dfrac{\text{OH}\times \text{aire(ABC)}}{3}$.

Ce volume vaut 2 et $\text{OH} =\dfrac{3\sqrt{22}}{11}$,
donc $2=\dfrac{\frac{3\sqrt{22}}{11} \times \text{aire(ABC)}}{3}$.

On en déduit que:
$\text{aire(ABC)} 
= \dfrac{6}{\text{OH}}
= \dfrac{6}{\frac{3\sqrt{22}}{11}}
= \dfrac{6\times 11}{3\sqrt{22}}
= \dfrac{22}{\sqrt{22}}
= \sqrt{22}$.

\item Soit la propriété: pour le tétraèdre OABC, \og le carré de l'aire du triangle ABC est égal à la somme des carrés des aires des trois autres faces du tétraèdre \fg.



\begin{list}{\textbullet}{}
\item $\text{aire(OAB)} = 3$ donc $\left (\text{aire(OAB)}\right )^2 = 9$

\item $\text{aire(OAC)} = \dfrac{\text{OA} \times \text{OC}}{2}=\dfrac{3\times 2}{2}=3$ donc $\left (\text{aire(OAC)}\right )^2 = 9$

\item $\text{aire(OBC)} = \dfrac{\text{OB} \times \text{OC}}{2}=\dfrac{2\times 2}{2}=2$  donc $\left (\text{aire(OBC)}\right )^2 = 4$
\end{list}
Donc $\left (\text{aire(OAB)}\right )^2 + \left (\text{aire(OAC)}\right )^2 + \left (\text{aire(OBC)}\right )^2 = 9+9+4=22$.

Comme $\text{aire(OABC)} = \sqrt{22}$, on a:\\
 $\left (\text{aire(OABC)}\right )^2=22=\left (\text{aire(OAB)}\right )^2 + \left (\text{aire(OAC)}\right )^2 + \left (\text{aire(OBC)}\right )^2$.

La propriété est vérifiée.
\end{enumerate}
